% Generated draft: requires mathematical and source review.
\documentclass[11pt]{article}
\usepackage[margin=25mm]{geometry}
\usepackage{fontspec}
\setmainfont{Latin Modern Roman}
\usepackage{amsmath,amssymb,mathtools,unicode-math}
\setmathfont{Latin Modern Math}
\usepackage{enumitem,needspace,xurl,hyperref}
\setlength{\parindent}{0pt}
\setlength{\parskip}{5pt}
\setlength{\emergencystretch}{3em}
\newcommand{\N}{\mathbb{N}}
\newcommand{\Z}{\mathbb{Z}}
\newcommand{\Q}{\mathbb{Q}}
\newcommand{\R}{\mathbb{R}}
\newcommand{\C}{\mathbb{C}}
\newcommand{\F}{\mathbb{F}}
\newcommand{\A}{\mathbb{A}}
\newcommand{\PP}{\mathbb{P}}
\newcommand{\E}{\mathbb{E}}
\newcommand{\eps}{\varepsilon}
\newcommand{\dd}{\,\mathrm{d}}
\begin{document}
\section*{UnsolvedMath ID: 800744}
\textbf{Source problem code:} catalogue entry
\subsection{Definitions and mathematical statement}
A Banach space is separable if it has a countable dense subset. Two normed spaces over the same scalar field are isomorphic, written \(Y\cong Z\), if there is a bounded linear bijection between them with bounded inverse. A Banach space is Hilbertian if it is isomorphic to a Hilbert space; non-Hilbertian means that no such isomorphism exists.

\textbf{Source question (G. Godefroy, reference label [Go]).} Let \(X\) be a separable non-Hilbertian Banach space. Does there exist a sequence \((X_n)_{n\geq 1}\) of subspaces of \(X\), endowed with their inherited norms, such that
\[
X_m\not\cong X_n\qquad\text{for every }m,n\geq 1\text{ with }m\neq n?
\]
The supplied source does not specify whether the scalar field is \(\R\) or \(\C\), whether the subspaces must be closed, or whether they must be infinite-dimensional. These qualifications remain unresolved. In particular, the intended admissibility of finite-dimensional subspaces must be clarified before this can be treated as a precise substantive open problem. Under the additional interpretation that all subspaces are closed and infinite-dimensional, the requested sequence would satisfy
\[
\forall n\geq 1,\quad X_n\subseteq X\text{ is a closed linear subspace and }\dim X_n=\infty,
\qquad
\forall m\neq n,\quad X_m\not\cong X_n.
\]
This additional interpretation is a conditional formulation, not an explicit hypothesis in the supplied excerpt.

The source introduces the question by reporting that every separable non-Hilbertian Banach space contains two non-isomorphic subspaces. That contextual assertion is retained as supplied background, with the same unresolved subspace conventions.
\subsection{Short English statement}
Godefroy asks whether every separable Banach space that is not isomorphic to a Hilbert space contains a countably infinite family of pairwise non-isomorphic subspaces. The source leaves the scalar field, closedness, and infinite-dimensionality requirements unstated; the intended subspace convention requires review.
\subsection{Sources}
\begin{itemize}
\item Local UnsolvedMath registry: permanent ID \(800744\), source problem code \texttt{catalogue entry}, title \emph{Banach Spaces and Ramsey Theory --- Problem 7}. The registry supplies the status \texttt{open}; this label is not an independently verified current mathematical assessment.
\item Dodos and Lopez-Abad, \emph{Banach Spaces and Ramsey Theory: Some Open Problems} (2010), Problem 7. Supplied source link: \url{https://arxiv.org/abs/1006.2668}. The record identifies the extraction method as \texttt{arxiv\_latex} and gives an access date of 2026-10-16. These are supplied provenance statements, not a claim of independent source access or verification.
\item The problem is explicitly attributed in the supplied statement to G. Godefroy, with reference label [Go]. The bibliographic entry corresponding to [Go] is not included in the supplied record. The registry lists Dodos and Lopez-Abad as proposers and 2010 as the proposed year.
\end{itemize}
\end{document}
